% ============ 第1章 §1.1 ℓ^p 空间 (原书 pp.1-6) ============
\section{\texorpdfstring{$\ell^{p}$ 空间}{ellp 空间}}\label{sec:1.1}
\index{0symbollinf@$\ell^{\infty}$}

\index{0symbollp@$\ell^{p}$}

\index{C@Cauchy--Schwarz 不等式}


设 $N$ 为自然数, 考虑由 $N$ 个实数组成的数列 $a_{1}, a_{2}, \ldots, a_{N}$. 将其视为 $N$ 维向量, 记 $a = (a_{1}, a_{2}, \ldots, a_{N})$. 这样的数列全体亦即向量全体构成线性空间 $\R^{N}$. 加法按分量定义, 数乘通过将该数乘到各分量来定义. 由数列构成的线性空间及其子空间称为数列空间.

\begin{definition}\label{def:1.1}
对 $1 \leq p < \infty$, 定义
\[
\| a \|_{\ell^{p}} = \left(\sum_{n=1}^{N} |a_{n}|^{p}\right)^{\frac{1}{p}},
\]
称之为 $a$ 的 $\ell^{p}$-范数. 另外, $\sum_{j=1}^{N} |a_{j}|^{p}$ 称为 $p$ 乘和.

当把线性空间 $\R^{N}$ 与 $\ell^{p}$-范数成对考虑时, 称 $\R^{N}$ 为有限集合 $\{1, 2, \ldots, N\}$ 上的 $\ell^{p}$ 空间.
\end{definition}

\begin{example}\label{exa:1.1}
当 $\{a_{n}\}_{n=1}^{N}$ 由 $a_{n} = n$ 给定时,
\[
\| \{a_{n}\}_{n=1}^{N} \|_{\ell^{1}} = 1 + 2 + \cdots + N = \frac{N(N+1)}{2},
\]
\[
\| \{a_{n}\}_{n=1}^{N} \|_{\ell^{2}} = \sqrt{1^{2} + 2^{2} + \cdots + N^{2}} = \sqrt{\frac{N(N+1)(2N+1)}{6}}.
\]
\end{example}

下面的不等式将在全书范围内说明 $\ell^{p}$-范数为何重要.

\begin{theorem}\label{thm:1.1}
对两个数列 $a = (a_{1}, a_{2}, \ldots, a_{N})$, $b = (b_{1}, b_{2}, \ldots, b_{N})$, 称为三角不等式的下述不等式成立:
\[
\|a + b\|_{\ell^{p}} \leq \|a\|_{\ell^{p}} + \|b\|_{\ell^{p}}.
\]
\end{theorem}

\begin{proof}
暂时假定 $a_{n}, b_{n} \geq 0$ 来考虑. 如果 $a$, $b$ 中某一方为 $0$, 则不等号变为等号. 基于这一点, 考虑所有分量都不为 $0$ 的两个数列 $a = (a_{1}, a_{2}, \ldots, a_{N})$, $b = (b_{1}, b_{2}, \ldots, b_{N})$. 这里使用三角不等式. 由函数 $\varphi(t) = t^{p}$ 的凸性, 即对 $t_{1}, t_{2} \geq 0$, $0 \leq \theta \leq 1$ 的不等式
\[
\varphi((1-\theta)t_{1} + \theta t_{2}) \leq (1-\theta)\varphi(t_{1}) + \theta\varphi(t_{2}),
\]
对 $\theta = \frac{\|b\|_{\ell^{p}}}{\|a\|_{\ell^{p}} + \|b\|_{\ell^{p}}}$, 有
\begin{align*}
\left(\frac{|a_{n} + b_{n}|}{\|a\|_{\ell^{p}} + \|b\|_{\ell^{p}}}\right)^{p} & \leq \left(\frac{|a_{n}| + |b_{n}|}{\|a\|_{\ell^{p}} + \|b\|_{\ell^{p}}}\right)^{p} \\
& \leq (1-\theta)\left(\frac{|a_{n}|}{\|a\|_{\ell^{p}}}\right)^{p} + \theta\left(\frac{|b_{n}|}{\|b\|_{\ell^{p}}}\right)^{p},
\end{align*}
所以, 将其对 $n = 1, 2, \ldots, N$ 求和并整理, 所期望的不等式成立.
\end{proof}

于是对任意数列 $a$, $b$ 有 $\|a + b\|_{\ell^{p}} \leq \|a\|_{\ell^{p}} + \|b\|_{\ell^{p}}$ 成立. 这个不等式称为 Minkowski 不等式. $p = 2$ 的情形就是通常的三角不等式. 接下来, 为了证明 Hölder 不等式, 先证明 Young 不等式.

\begin{lemma}[Young 不等式]\label{lem:1.1}
设 $p, q > 1$ 满足 $\frac{1}{p} + \frac{1}{q} = 1$. 此时, 对 $\alpha, A, B > 0$,
\[
\frac{\alpha^{p}}{p}A + \frac{\alpha^{-q}}{q}B \geq A^{\frac{1}{p}}B^{\frac{1}{q}}
\]
成立. 当 $\alpha = \left(\frac{B}{A}\right)^{\frac{1}{p+q}}$ 时等号成立. 特别地, 对 $a, b > 0$,
\[
ab \leq \frac{1}{p}a^{p} + \frac{1}{q}b^{q}
\]
成立.
\end{lemma}

\begin{proof}
令 $f(\alpha) = \frac{\alpha^{p}}{p}A + \frac{\alpha^{-q}}{q}B$, 则 $f'(\alpha) = \alpha^{p-1}A - \alpha^{-q-1}B$. 因此, 当 $\alpha = \left(\frac{B}{A}\right)^{\frac{1}{p+q}}$ 时 $f$ 取极小值, 所以
\begin{align*}
f(\alpha) & \geq f\left(\left(\frac{B}{A}\right)^{\frac{1}{p+q}}\right) \\
& = \frac{1}{p}\left(\frac{B}{A}\right)^{\frac{p}{p+q}}A + \frac{1}{q}\left(\frac{B}{A}\right)^{-\frac{q}{p+q}}B \\
& = \frac{1}{p}A^{\frac{q}{p+q}}B^{\frac{p}{p+q}} + \frac{1}{q}A^{\frac{q}{p+q}}B^{\frac{p}{p+q}} \\
& = A^{\frac{1}{p}}B^{\frac{1}{q}}.
\end{align*}
得证.
\end{proof}

应用 Young 不等式可知如下事实. 设 $A, B, C > 0$. 又设 $p, q > 1$ 满足 $\frac{1}{p} + \frac{1}{q} = 1$. 此时, 若对 $\alpha > 0$,
\[
\frac{\alpha^{p}}{p}A + \frac{\alpha^{-q}}{q}B \geq C
\]
成立, 则 $C \leq A^{\frac{1}{p}}B^{\frac{1}{q}}$ 成立. 这个想法在本节剩余部分将频繁使用. 在引理~\ref{lem:1.1} 中, 由关系式 $\frac{1}{p} + \frac{1}{q} = 1$ 联系起来的 $1 < p, q < \infty$ 称为处于调和共轭关系. 进一步, 对 $1 < p < \infty$, 称 $q = \frac{p}{p-1}$ 为 $p$ 的调和共轭, 也常记为 $p'$. 另外, 约定 $1' = \infty$, $\infty' = 1$.

\begin{example}\label{exa:1.2}
$1$ 的调和共轭是 $\infty$, $2$ 的调和共轭是 $2$, $3$ 的调和共轭是 $\frac{3}{2}$.
\end{example}

对于处于调和共轭关系的 $p$, $q$, 联系 $\ell^{1}$, $\ell^{p}$, $\ell^{q}$ 的 Hölder 不等式是重要的.

\begin{theorem}[Hölder 不等式]\label{thm:1.2}
设 $a_{1}, a_{2}, \ldots, a_{N}, b_{1}, b_{2}, \ldots, b_{N}$ 为实数. 又对处于调和共轭关系的 $1 < p, q < \infty$, 不等式
\[
\left| \sum_{n=1}^{N} a_{n}b_{n} \right| \leq \left(\sum_{n=1}^{N} |a_{n}|^{p}\right)^{\frac{1}{p}}\left(\sum_{n=1}^{N} |b_{n}|^{q}\right)^{\frac{1}{q}}
\]
成立.
\end{theorem}

$p = q = 2$ 的情形称为 Cauchy--Schwarz 不等式.

\begin{proof}
必要时将 $a_{j}$ 替换为 $|a_{j}|$, 将 $b_{j}$ 替换为 $|b_{j}|$ 来考虑, 可设 $a_{j}, b_{j} \geq 0$. 实际上, 由三角不等式,
\[
\left| \sum_{n=1}^{N} a_{n}b_{n} \right| \leq \sum_{n=1}^{N} |a_{n}b_{n}|
\]
就是这个缘故. 以 $\alpha > 0$ 为常数, 由引理~\ref{lem:1.1} 的不等式, 令 $a = \alpha a_{n}$, $b = \alpha^{-1}b_{n}$, 得
\begin{equation}\tag{1.1}\label{eq:1.1}
a_{n}b_{n} \leq \frac{\alpha^{p}}{p}a_{n}^{p} + \frac{\alpha^{-q}}{q}b_{n}^{q}
\end{equation}
将其对 $n = 1, 2, \ldots, N$ 求和.
\begin{equation}\tag{1.2}\label{eq:1.2}
\sum_{n=1}^{N} a_{n}b_{n} \leq \frac{\alpha^{p}}{p}\sum_{n=1}^{N} a_{n}^{p} + \frac{\alpha^{-q}}{q}\sum_{n=1}^{N} b_{n}^{q}.
\end{equation}
这里, 将右边的式子看作 $\alpha$ 的函数, 令
\[
g(\alpha) = \frac{\alpha^{p}}{p}\sum_{n=1}^{N} a_{n}^{p} + \frac{\alpha^{-q}}{q}\sum_{n=1}^{N} b_{n}^{q}.
\]
由引理~\ref{lem:1.1}, $g(\alpha) \geq \sum_{n=1}^{N} a_{n}b_{n}$, 所以利用 $g(\alpha)$ 的最小值,
\[
\sum_{n=1}^{N} a_{n}b_{n} \leq \left(\sum_{n=1}^{N} a_{n}^{p}\right)^{\frac{1}{p}}\left(\sum_{n=1}^{N} b_{n}^{q}\right)^{\frac{1}{q}}
\]
得证.
\end{proof}

来考察不等式的等号成立条件吧. 虽说并非所有不等式的等号成立条件都能考察, 但在能够考察等号成立条件的时候, 考察一下是很重要的.

考虑最简单的最初的定理的情形. 由于式~\eqref{eq:1.1} 中等号必须成立, 所以参考 $f(a) = \frac{1}{p}a^{p} + \frac{1}{q}b^{q} - ab$ 的极小值在哪里取到,
\begin{equation}\tag{1.3}\label{eq:1.3}
\alpha a_{n} = \left(\alpha^{-1}b_{n}\right)^{\frac{1}{p-1}}, \quad n = 1, 2, \ldots, N
\end{equation}
必须成立. 这里, 式~\eqref{eq:1.2} 的右边成为 Hölder 不等式的右边, 所以 $\alpha$ 取使式~\eqref{eq:1.2} 的右边最小的值.

反过来, 若条件 ``存在某个 $\beta$ 使得 $a_{n} = \beta(b_{n})^{\frac{1}{p-1}}, n = 1, 2, \ldots, N$ 成立'' 成立, 则令 $\beta = \alpha^{-\frac{p}{p-1}}$, 关系式~\eqref{eq:1.3} 就出现了. 利用它可知 Hölder 不等式的等号成立. 即使 $a_{n}$, $b_{n}$ 是一般的实数, 只要仔细考虑符号, 情况也是相同的. 也就是说, $a_{n}b_{n} \geq 0$ 对所有 $n = 1, 2, \ldots, N$ 成立, 或 $a_{n}b_{n} \leq 0$ 对所有 $n = 1, 2, \ldots, N$ 成立, 并且存在某个 $\beta$ 使得 $|a_{n}| = \beta|b_{n}|^{\frac{1}{p-1}}, n = 1, 2, \ldots, N$ 成立, 这种情形就是等号成立条件.

\begin{definition}\label{def:1.2}
对数列 $a = \{a_{n}\}_{n=1}^{N}$, 定义 $\|a\|_{\ell^{\infty}} = \sup_{n = 1, 2, \ldots, N} |a_{n}|$.
\end{definition}

有时也给 $\R^{N}$ 配备 $\ell^{\infty}$-范数.

可以看出, 到目前为止的命题对 $p = 1, \infty$ 也成立.

数列的 Hölder 不等式还可以变形为积分形式. 设 $f, g: [a, b] \to \R$ 为连续有界函数. 由引理~\ref{lem:1.1}, 对任意 $A > 0$,
\[
|f(x)g(x)| \leq \frac{A^{p}}{p}|f(x)|^{p} + \frac{A^{-q}}{q}|g(x)|^{q}
\]
成立. 将其在 $[a, b]$ 上积分,
\[
\int_{a}^{b} |f(x)g(x)| \dd x \leq \frac{A^{p}}{p}\int_{a}^{b} |f(x)|^{p} \dd x + \frac{A^{-q}}{q}\int_{a}^{b} |g(x)|^{q} \dd x
\]
成立. 由引理~\ref{lem:1.1}, 选取 $A$ 使右边最小, 得
\begin{equation}\tag{1.4}\label{eq:1.4}
\int_{a}^{b} |f(x)g(x)| \dd x \leq \left(\int_{a}^{b} |f(x)|^{p} \dd x\right)^{\frac{1}{p}}\left(\int_{a}^{b} |g(x)|^{q} \dd x\right)^{\frac{1}{q}}.
\end{equation}
通过替换幂次, 对满足 $\frac{1}{r} = \frac{1}{p} + \frac{1}{q}$ 的 $p, q, r \geq 1$, 得
\begin{equation}\tag{1.5}\label{eq:1.5}
\left(\int_{a}^{b} |f(x)g(x)|^{r} \dd x\right)^{\frac{1}{r}} \leq \left(\int_{a}^{b} |f(x)|^{p} \dd x\right)^{\frac{1}{p}}\left(\int_{a}^{b} |g(x)|^{q} \dd x\right)^{\frac{1}{q}}.
\end{equation}
式~\eqref{eq:1.4} 与式~\eqref{eq:1.5} 称为 Hölder 不等式的积分形式. 即使在积分的情形, $p = q = 2$ 时也称为 Cauchy--Schwarz 不等式. 例如, 在 Hölder 不等式~\eqref{eq:1.5} 中, 令 $p = 3$, $q = 6$, $r = 2$, $f = |\varphi|^{\frac{1}{3}}$, $g = |\varphi|^{\frac{2}{3}}$, 则对函数 $\varphi: [0, 1] \to \R$ 得积分不等式
\begin{equation}\tag{1.6}\label{eq:1.6}
\left(\int_{0}^{1} |\varphi(x)|^{2} \dd x\right)^{\frac{1}{2}} \leq \left(\int_{0}^{1} |\varphi(x)| \dd x\right)^{\frac{1}{3}}\left(\int_{0}^{1} |\varphi(x)|^{4} \dd x\right)^{\frac{1}{6}}.
\end{equation}
Hölder 不等式的积分形式不限于连续函数, 只要积分能用 Riemann 和的极限表示就足够了, 因此例如对只在有限个点不连续的有界函数等也成立. 进一步, 例如对这种函数与 $1 \leq p \leq p_{0} < \infty$,
\begin{equation}\tag{1.7}\label{eq:1.7}
\left(\int_{0}^{1} |f(x)|^{p} \dd x\right)^{\frac{1}{p}} \leq \left(\int_{0}^{1} |f(x)|^{p_{0}} \dd x\right)^{\frac{1}{p_{0}}}
\end{equation}
成立. 实际上, 在式~\eqref{eq:1.4} 中取 $[a, b] = [0, 1]$, $g = 1$, 则
\[
\int_{0}^{1} |f(x)| \dd x \leq \left(\int_{0}^{1} |f(x)|^{p} \dd x\right)^{\frac{1}{p}},
\]
然后适当替换函数和指数即可.

\begin{exercise}\label{exer:1.1}
设数列 $a = \{a_k\}_{k=1}^{N}$ 由 $a_k = \sqrt{k} \ (k = 1, 2, \ldots, N)$ 给出. 此时, 计算 $\|a\|_{\ell^2}$, $\|a\|_{\ell^4}$.
\end{exercise}

\begin{exercise}\label{exer:1.2}
设 $\vec{a} = (2, 3, 1)$, $\vec{b} = (3, 4, 2)$. 计算 $\vec{a} + \vec{b}$, 并利用其结果证明 $\sqrt{83} \leq \sqrt{14} + \sqrt{29}$.
\end{exercise}

\begin{exercise}\label{exer:1.3}
对 4 个正实数 $a, b, c, d$, 证明不等式 $(ac^{2} + bd^{2})^{3} \leq (a^{3} + b^{3})(c^{3} + d^{3})^{2}$.
\end{exercise}

\begin{exercise}\label{exer:1.4}
对 $x, y \in \R$ 与 $\theta \in (0, 1)$, 证明 $e^{(1-\theta)x + \theta y} \leq (1-\theta)e^{x} + \theta e^{y}$.
\end{exercise}

\begin{exercise}\label{exer:1.5}
对 $1 \leq p < \infty$, $k \in \R$ 与不为零数列的实数列 $a = \{a_k\}_{k=1}^{N}$, 计算 $\|ka\|_{\ell^{p}} \div \|a\|_{\ell^{p}}$.
\end{exercise}

\begin{exercise}\label{exer:1.6}
设 $a, b \in \R$, 证明
\[
\max(a, b) = \frac{a + b + |a - b|}{2}, \quad \min(a, b) = \frac{a + b - |a - b|}{2}.
\]
\end{exercise}
